\section{Complete cancellation collections for fourteen values}\label{sec:linear}
We now return to linear forms. The target set is
$\{\beta(4),\beta(6),\ldots,\beta(30)\}$, with fourteen members.
Fifteen independent forms on these constants together with $1$ cannot
all vanish. This collective nonvanishing route avoids having to identify
one particular form that remains nonzero. Its unresolved cost is the
height of the entire collection.

The rational-function construction adapts Krisman's hypergeometric
method \cite{Krisman} and the isolated-pole refinement in
\cite{fifteen}. We use twenty-nine nested pole blocks, double integer
zeros, and one derivative. These choices differ from the thirty-one
blocks, triple zeros, and two derivatives in the fifteen-value result.

\subsection{Formal rank from a rational antidifference}
It is convenient to prove rank for general nested blocks. Let $n>0$
be even, let $b$ be a nonnegative integer, and let $S$ be a positive
odd integer. Choose integers $0\leq d_j\leq n/2$ for $1\leq j\leq S$.
Put $L=n+2b$ and
\[
 N(t)=\prod_{\ell=-n-b-1}^{b}(t-\ell)^2,\quad
 Q(t)=\prod_{j=1}^S\prod_{a=d_j}^{n-d_j}(t+a+\tfrac12),
 \quad R(t)=N(t)/Q(t).
\]
For a polynomial invariant under $t\mapsto-n-1-t$, with $PR$ proper,
define
\[
 \mathcal J(P)=\sum_{t\geq0}(-1)^t(PR)'(t).
\]
The differentiated series converges absolutely. Partial fractions
give rational formal coefficients on $1,\beta(2),\ldots,\beta(S+1)$.
A \emph{formal zero} means that all these coefficients vanish, without
making any assumption about relations among the actual constants.

\begin{theorem}\label{thm:rank}
If $\deg P\leq S+1$, a formal zero forces $P=0$. More generally,
subject to properness, all formal-zero symmetric polynomials are
\begin{equation}\label{eq:telescopers}
 P(t)=g(t)B(t-1)+g(t+1)A(t),\qquad g(-n-t)=-g(t),
\end{equation}
where
\[
 A(t)=(t+n+b+2)^2\prod_j(t+d_j+\tfrac12),\quad
 B(t)=(t-b)^2\prod_j(t+n-d_j+\tfrac32).
\]
The first possible nonzero formal identity has degree $S+3$.
\end{theorem}
\begin{proof}
Write $F=PR=\sum_{a,i}c_{i,a}(t+a+1/2)^{-i}$, with pole orders
$B_a=\#\{j:d_j\leq a\leq n-d_j\}$. Reflection gives
$c_{i,n-a}=(-1)^{i+1}c_{i,a}$. The coefficient of $\beta(i+1)$ is
$-i2^{i+1}\sum_a(-1)^ac_{i,a}$. Reflection already makes this sum
zero for even $i$. If the other beta coefficients vanish too, all
alternating residue sums vanish.

Define $h_{i,-1}=0$, $h_{i,a}=c_{i,a}-h_{i,a-1}$. Then $h_{i,n}=0$,
and matching principal parts gives the unique proper rational function
\[
 H(t)=\sum_{a=0}^{n-1}\sum_i\frac{h_{i,a}}{(t+a+1/2)^i},
 \qquad F(t)=H(t)+H(t+1).
\]
Uniqueness follows because a nonzero rational solution of
$H(t)+H(t+1)=0$ would have infinitely many translated poles, or would
be a polynomial whose leading terms cannot cancel.

The forward and backward recurrences bound the pole order at $-a-1/2$
by $\min(B_a,B_{a+1})$, since the orders are unimodal. Thus
\[
 H=V/Q_H,\quad
 Q_H=\prod_j\prod_{a=d_j}^{n-d_j-1}(t+a+1/2),\quad
 \deg Q_H=\deg Q-S.
\]
This loss of one factor per block is the rank mechanism. A nonzero
$H$ has the same degree at infinity as $F$, because its leading term
doubles when shifted and added. Consequently
$\deg V=2L+4+\deg P-S$.

Uniqueness and reflection imply $H(-n-t)=-H(t)$, so $H(-n/2)=0$.
The double zeros of $F$ imply the two recurrences
$H(\ell)+H(\ell+1)=0$ and $H'(\ell)+H'(\ell+1)=0$ on the
integer-zero interval. The first propagates from the center to every
integer $-n-b-1,\ldots,b+1$. The second needs the rational constant:
telescoping gives $\mathcal J(P)=H'(0)$, so a formal zero gives
$H'(0)=0$ and the same propagation for derivatives. Thus $V$ has
$L+3$ double zeros, total degree at least $2L+6$. If $\deg P\leq S+1$,
its degree is at most $2L+5$, a contradiction unless $P=0$.

In general those double zeros give
$H=g(t)(t-b-1)^2N(t)/Q_H(t)=g(t)B(t-1)R(t)$.
Reflection makes $g$ odd about $-n/2$, and
$R(t+1)/R(t)=A(t)/B(t)$ gives \eqref{eq:telescopers}. Conversely that
formula yields $PR=H+H(t+1)$; properness makes $H$ proper, the
alternating residue sums vanish, and $H'(0)=0$ because of its double
zero there. Hence every formal coefficient vanishes. Finally the
leading terms in \eqref{eq:telescopers} add, giving
$\deg P=\deg g+S+2\geq S+3$.
\end{proof}

\subsection{The sixteen coefficient directions and integral reduction}
Specialize to $116\mid n$, $b=n$, $S=29$, and offsets
$d_j=jn/116$ for $0\leq j\leq28$. Then $\deg N=6n+4$ and
$\deg Q=22n+29$. Put $x=2t+n+1$. The sixteen multipliers
$1,x^2,\ldots,x^{30}$ give an invertible rational coefficient matrix
on $1,\beta(2),\ldots,\beta(30)$ by Theorem~\ref{thm:rank}.
The beta-two cancellation space has dimension fifteen, and every
independent collection of fifteen cancellations gives an invertible
matrix on $1,\beta(4),\ldots,\beta(30)$.

Higher degrees add known identities rather than new coefficient
directions. Write $X_n=3n+3$ and let $E_h(x)$ be the even part of
\[
 (x-1)^{2h+1}(x-X_n)^2\prod_{j=0}^{28}(x+n-2d_j),\qquad h\geq0.
\]
Taking an even part just discards odd coefficients. Thus $E_h$ is
monic in $\Z[x^2]$, of degree $h+16$ in $x^2$. Formula
\eqref{eq:telescopers} shows that it is a formal zero whenever the
product is proper. Subtracting integer multiples of these monic
identities reduces every integer even polynomial of degree $2D\leq2n$
uniquely to degree at most thirty, with the same coefficient vector.
This is an additive quotient; multiplying an arbitrary identity by
$x^2$ need not preserve it.

For the weighted coefficient norm $\|P\|_{X_n}=\sum_k|[x^k]P|X_n^k$,
\[
 \|E_h\|_{X_n}/X_n^{2h+32}
 \leq4(1+1/X_n)^{2h+1}(1+n/X_n)^{29}<40000
\]
in this range. Successive monic elimination therefore gives
$\|P_{\rm red}\|_{X_n}\leq40001^{\max(0,D-15)}\|P\|_{X_n}$.
For $D=o(n)$ this costs only $\exp(o(n))$, and introduces no
denominators. It does not ensure that short vectors have independent
nonzero reductions.

\subsection{Eleven directions with uniform height and decay}
For $P(x)=\sum_{j=0}^{15}w_jx^{2j}$ put
$h_n(P)=\sum_j|w_j|(3n+3)^{2j}$, and let $\Lambda_n\subset\Z^{16}$
be the lattice canceling the beta-two coefficient. Its rank is fifteen.
The following two quantitative estimates use different clearing integers:
one bounds the primitive normal to this lattice, the other clears
the actual forms.

\begin{lemma}\label{lem:linear-inputs}
Let $a_n$ be the primitive integer normal to $\Lambda_n$. There is
a constant $H_*<24.522$ such that
\[
 \log\|a_n\|_\infty\leq(H_*+o(1))n.
\]
There are common positive integers $D_n^*$ clearing every coefficient
of every allowed beta-two-canceling form, for which
\[
 |D_n^*\mathcal J_n(P)|
 \leq h_n(P)\exp((-C_*+o(1))n),\qquad C_*>6.0384.
\]
The error terms are uniform in the multiplier coefficients.
\end{lemma}
Appendix~\ref{app:linear-inputs} gives the arithmetic and analytic
estimates, including the finite rational certificate defining these
constants. The parity grouping used for $a_n$ is especially useful:
the even-address and odd-address denominator products are coprime
modulo two, so their grouped residue sums have much better dyadic
integrality than individual residues.

\begin{theorem}\label{thm:eleven}
For every sufficiently large positive multiple $n$ of $116$, there
is an integer matrix $C_n$ with eleven rows and fifteen columns such
that, for $\theta=(1,\beta(4),\beta(6),\ldots,\beta(30))^T$,
\[
 \rank_{\Q}C_n=11,\qquad \|C_n\theta\|_\infty\leq e^{-1.12n}.
\]
Its rows come from eleven independent beta-two-canceling multipliers
of weighted height at most $e^{4.91n}$.
\end{theorem}
\begin{proof}
The lattice with primitive normal $a_n$ has Euclidean covolume
$\|a_n\|_2$. After applying
$T=\diag(1,X_n^2,\ldots,X_n^{30})$, its covolume is
\[
 \det T\,\|T^{-1}a_n\|_2\leq4X_n^{240}\|a_n\|_\infty.
\]
The classical LLL product estimate \cite[Proposition~1.6]{LLL} gives
a basis whose product of lengths is at most
$2^{15\cdot14/4}$ times this covolume. To recall the relevant estimate,
size reduction gives $|\mu_{ij}|\leq1/2$ and the reduction condition
gives $\|b_i^*\|^2\geq\|b_{i-1}^*\|^2/2$. Expanding the basis
vectors in their orthogonal components and multiplying yields the
stated product bound.

Order its lengths as $s_1\leq\cdots\leq s_{15}$. Each is at least
one since the weighted lattice is integral, and
$s_k^{16-k}\leq\prod_i s_i$. Weighted $\ell^1$ height is at most
four times Euclidean length. Thus the $k$th weighted successive
minimum satisfies
\[
 \limsup n^{-1}\log\lambda_k(n)\leq H_* /(16-k).
\]
At $k=11$ this is below $4.905$, so eleven independent vectors have
height at most $e^{4.91n}$ for all sufficiently large admissible $n$.
Apply the \emph{same} $D_n^*$ from Lemma~\ref{lem:linear-inputs} to
their forms. Since $6.0384-4.91>1.12$, their values satisfy the claimed
bound. The invertible formal coefficient map preserves their
independence after cancellation and common clearing.
\end{proof}

An eleven-by-fifteen matrix can annihilate $\theta$ exactly, so
Theorem~\ref{thm:eleven} does not establish irrationality. Fifteen
independent small rows would exclude that possibility. The product
argument above gives only the twelfth-minimum upper rate
$H_*/4=6.13044\ldots$, already above $C_*$. This is the limitation
of this height estimate, not a lower bound on the true twelfth or
fifteenth minimum.

At $n=116,232,464$, exact finite computations do give complete
fifteen-row collections, with largest primitive values below
$10^{-236}$, $3.68\times10^{-470}$, and $6.04\times10^{-890}$,
respectively. These are the same kind of objects needed for the
contradiction. The unproved step is their uniform height or decay
along an unbounded sequence. A bound for one shortest vector is
insufficient, and the available covolume bound does not give the
required rate: the kernel of
$(1,M,0,\ldots,0)$ has many unit directions while its last direction
must still have size comparable to $M$.
